\section{Related work}
\subsection{LLM and VLM as a judge}


LLM-as-a-judge can provide scalable reward signals for RL fine-tuning, reducing dependence on expensive human annotations while enabling more efficient preference learning. In text-only domains, large language models (LLMs) serve as automatic evaluators for tasks including summarization, dialogue, and reasoning, often achieving high consistency with human annotators.Recent work has explored various training approaches for LLM-as-a-judge systems~\cite{zheng2023llmjudge,lee2023rlaif,yuan2024srlm,cui2023ultrafeedback,lambert2024tulu, whitehouse2025j1incentivizingthinkingllmasajudge, saha2025learningplanreason}.

Similarly, vision-language models (VLMs) have been deployed as multimodal judges for captioning, visual question answering (VQA), and reasoning tasks. 
Training approaches for model evaluators have primarily followed two directions: scaling human preference collection and distilling knowledge from large closed-source models such as GPT and Claude (which themselves rely heavily on human annotation)~\cite{llava-critic, wang2025unified, vlrb, mmrb}. Some works focus on building specialized judges for single dimensions like hallucination detection using synthetic data, while general-purpose judges typically require substantial amounts of preference data~\cite{jing2024fgaif, 2023llavarlhf}.

\subsection{Synthetic Data for Model Self-Improvement}
Synthetic data has emerged as a central driver of self-improvement in both LLMs and VLMs, enabling models to generate, evaluate, and refine their own outputs. Instruction-tuning pipelines bootstrap new tasks from model-generated data to expand task coverage without human labels, while iterative refinement frameworks enable models to critique and improve their own responses~\cite{Wang2022SelfInstruct,Li2024GLAN,Madaan2023SelfRefine,Shinn2023Reflexion}. Self-rewarding and self-play approaches further demonstrate that models can generate training signals to surpass supervised baselines~\cite{alemohammad2024self}. However, uncontrolled synthetic data generation can lead to degenerative feedback loops and preference leakage, underscoring the importance of careful curation~\cite{Alemohammad2023MAD,Li2025PreferenceLeakage,liu2024bestpracticeslessonslearned}.

The multimodal domain presents greater opportunities and unique challenges. While scarce and expensive multimodal annotations create stronger incentives for synthetic data generation, VLMs face distinctive issues including hallucination~\cite{sun2023aligninglargemultimodalmodels, povid, zhou2024analyzingmitigatingobjecthallucination,leng2023mitigatingobjecthallucinationslarge} and text-image misalignment~\cite{zhou2024aligningmodalitiesvisionlarge}. Most such work relies on larger models to create synthetic data. Our work extends text-domain self-improvement methods to address multimodal-specific challenges of maintaining visual grounding and factual accuracy without human annotations or larger models.